\chapter*{TÓM TẮT}
\vspace{-40pt}

Tóm tắt luận án cần được viết một cách ngắn gọn, súc tích nhưng phải phản ánh đầy đủ các nội dung cốt lõi của toàn bộ công trình nghiên cứu. Cấu trúc lý tưởng của một phần tóm tắt bao gồm các phần sau:

\begin{itemize}
    \item \textbf{Đặt vấn đề:} Nêu rõ lý do thực hiện đề tài, tầm quan trọng và mục tiêu chính của luận án (1-2 câu).
    \item \textbf{Phương pháp nghiên cứu:} Tóm tắt ngắn gọn cách tiếp cận, phương pháp luận, mô hình toán học, dữ liệu hoặc các thực nghiệm đã được sử dụng để giải quyết vấn đề (2-3 câu).
    \item \textbf{Kết quả nghiên cứu:} Trình bày những phát hiện, đóng góp mới hoặc kết quả nổi bật nhất mà luận án đã đạt được (2-3 câu).
    \item \textbf{Kết luận và Ý nghĩa:} Khẳng định lại giá trị khoa học, ý nghĩa thực tiễn hoặc hướng phát triển tiếp theo của đề tài (1-2 câu).
\end{itemize}

Lưu ý: Trong phần tóm tắt, hạn chế tối đa việc sử dụng các ký hiệu toán học quá phức tạp, viết tắt chưa qua định nghĩa, hoặc trích dẫn tài liệu tham khảo.

\vspace{15pt}
\noindent \textbf{\textit{Từ khóa:}} Từ khóa 1, Từ khóa 2, Từ khóa 3, Từ khóa 4 (Chọn từ 3 đến 5 cụm từ chuyên ngành đặc trưng nhất cho luận án).


